\documentclass[aspectratio=169,11pt]{beamer}
\usetheme{default}
\usecolortheme{dove}
\setbeamertemplate{navigation symbols}{}
\setbeamertemplate{footline}[frame number]
\usepackage[utf8]{inputenc}
\usepackage{graphicx}
\graphicspath{{../report/figures/}}   % reuse the figures exported by R

\title{How Big Are Prague's Cafés?}
\subtitle{Café Size --- Data Analytics Semestral Project}
\author{Anthony George \and \textit{teammates}}
\date{November 2026}

% One slide per finding: key plot + estimate with CI + effect size + plain-words meaning + test name.
% Present findings, not procedure. No assumption details here (they are in the report).
\newcommand{\finding}[4]{%
\begin{frame}{#1}
\begin{columns}[T]
\column{.58\textwidth}\centering\includegraphics[width=\linewidth]{#2}
\column{.4\textwidth}\small #3\par\vspace{1em}\footnotesize\textit{Test: #4}
\end{columns}
\end{frame}}

\begin{document}
\frame{\titlepage}

\begin{frame}{Introduction}
\begin{itemize}
  \item 195 Prague cafés visited in person by the class (autumn 2026)
  \item Our chapter: \textbf{how big are cafés}, where is outdoor seating, and does size go with ratings and prices?
  \item \textit{TODO: one line on how the team divides the ten findings}
\end{itemize}
\end{frame}

\finding{A1: A typical café seats about 28 inside}{A1_figure.pdf}
{Mean 28.0 seats, 95\% CI [24.3, 32.4], n = 193 --- no evidence it differs from 30. Median 22: most cafés are smaller, a few very large ones pull the mean up.}
{one-sample t-test, p = .331}

\finding{A2: Cafés with outdoor seating are a bit bigger inside}{A2_figure.pdf}
{+7.5 seats, 95\% CI [$-$0.4, 15.4], d = 0.26. \textit{TODO: plain-words meaning.}}
{Welch t-test, one-tailed p = .031}

\finding{A3: About 17 more seats inside than outside}{A3_figure.pdf}
{+17.3 seats, 95\% CI [13.5, 21.2], dz = 0.64. \textit{TODO}}
{paired t-test, p < .001}

\finding{A4: Outdoor seating differs between districts}{A4_figure.pdf}
{\textit{TODO: which districts differ (Dunn--Holm table).}}
{Kruskal--Wallis H(5) = 13.00, p = .023}

\finding{A5: Small cafés are more often top-rated}{A5_figure.pdf}
{58.5\% vs 29.7\%: +28.8 pp, 95\% CI [13.5, 44.1], V = 0.27. \textit{TODO}}
{$\chi^2$ test, one-tailed p < .001}

\finding{A6: Bigger cafés are slightly pricier}{A6_figure.pdf}
{Spearman $\rho$ = 0.25, 95\% CI [0.09, 0.39], n = 191. \textit{TODO}}
{Spearman correlation, p < .001}

\finding{A7: +0.7 CZK Price Level per 10 seats}{A7_figure.pdf}
{0.69 CZK per 10 seats, 95\% CI [0.07, 1.31], R$^2$ = 0.02. \textit{TODO}}
{simple regression, p = .030}

\finding{A8: Controlling for distance from the Orloj}{A8_figure.pdf}
{0.61 CZK per 10 seats, 95\% CI [$-$0.01, 1.23]; most of the slope remains but it is no longer clearly different from zero. \textit{TODO}}
{multiple regression, p = .055}

\finding{A9: The size--price link depends on outdoor seating}{A9_figure.pdf}
{Slope difference $-$2.62 CZK per 10 seats, 95\% CI [$-$4.15, $-$1.10]. \textit{TODO}}
{regression with interaction, p < .001}

\finding{A10: Chance of above-median prices per 10 seats}{A10_figure.pdf}
{+2.7 pp per 10 seats, 95\% CI [$-$0.3, 5.7] --- no clear evidence. \textit{TODO}}
{logistic regression, average marginal effect, p = .077}

\begin{frame}{Conclusion}
\begin{itemize}
  \item \textit{TODO: 3 bullets answering ``how big are Prague cafés, and does size matter?''}
  \item Limits: convenience sample, observational data, cafés in a district not independent
\end{itemize}
\end{frame}
\end{document}
